\documentclass[11pt,a4paper]{article}
\usepackage{fontspec}
\setmainfont{DejaVu Sans}[Scale=0.90]
\usepackage{amsmath,amssymb}
\usepackage[margin=2.2cm]{geometry}
\usepackage{booktabs,array,tabularx}
\usepackage{graphicx}
\graphicspath{{../figures/}}
\usepackage{enumitem}
\usepackage{titlesec}
\usepackage{parskip}
\titlelabel{\thetitle.\quad}
\titleformat{\section}{\normalfont\large\bfseries}{\thesection.}{0.6em}{}
\titleformat{\subsection}{\normalfont\normalsize\bfseries}{\thesubsection}{0.6em}{}
\titlespacing*{\section}{0pt}{15pt}{5pt}
\titlespacing*{\subsection}{0pt}{11pt}{4pt}
\setlist[itemize]{leftmargin=1.4em,itemsep=1pt,topsep=3pt}
% Preserve fonts/margins; permit a final line-breaking pass for long prose.
\setlength{\emergencystretch}{2em}
\begin{document}
\begin{center}
{\Large\bfseries Premier export microscopique de l'AlN}\\[4pt]
{\normalsize Comptage spectral, contrôles et limite de validation}\\[2pt]
{\small Partie I --- 25 septembre 2026}
\end{center}
\section{Objet et statut}
La note 19 fournit une référence finie exactement résonante. Cette note
examine un export obtenu à partir de constantes de force publiques
de l'AlN. Elle ne construit pas encore l'opérateur physique complet.
Le calcul emploie phono3py et phonopy 4.5.0, une grille $3^3$, $T=300$ K,
et des points externes sélectionnés. Les entrées et versions sont archivées
avec leurs empreintes dans \texttt{theory/aln/interaction\_pilot/}.

Les constantes harmoniques proviennent d'une supercellule $5\times5\times3$,
les constantes cubiques d'une supercellule $3\times3\times2$.
La correction polaire utilise le fichier BORN. Ces données partagées
constituent une hypothèse amont : leur exactitude physique n'est pas
prouvée par la reconstruction des sorties du même calcul.

\section{Grandeurs exportées}
La grille contient 27 points réguliers, représentés par 39 entrées de zone
de Brillouin avec duplications de frontière, et 12 branches.
Pour le premier point externe, huit triplets réduits portent les poids
\[
(2,1,4,2,4,2,8,4),\qquad \sum_t w_t=27.
\]
Une énumération indépendante des orbites par rotations entières et échange
des partenaires retrouve ces huit classes. Elle vérifie le comptage des
points pour la somme scalaire, pas les permutations complètes des branches
dans les sous-espaces dégénérés.

L'export \texttt{pp} contient les interactions au carré en eV$^2$.
Dans la conversion par défaut inspectée, le facteur $1/N_q$ est déjà
inclus. Le diviser encore par 27 serait incorrect.
Les fréquences et demi-largeurs sont exprimées en THz ordinaires.

\section{Reconstruction indépendante de la somme}
Avec $g_\sigma$ la gaussienne normalisée en fréquence et
$n_j=(\exp(hf_j/k_BT)-1)^{-1}$, la formule implémentée directement est
\[
\begin{split}
\gamma_0=C_\gamma\sum_{t,b_1,b_2}w_t\,P_{t,0,1,2}
\bigl\{&(n_1+n_2+1)g_\sigma(f_0-f_1-f_2)\\
&+(n_1-n_2)[g_\sigma(f_0+f_1-f_2)
-g_\sigma(f_0-f_1+f_2)]\bigr\}.
\end{split}
\]
Ici $P=\texttt{pp}$ et
\[
 C_\gamma=\frac{18\pi}{(\hbar\,2\pi\,10^{12})^2}
 =3\,306\,215.6970
\]
dans les unités mixtes eV et THz employées. Le seuil des fréquences
filles est $10^{-4}$ THz, conformément au défaut inspecté.
La convention de durée de vie est
$\tau^{-1}=4\pi\,10^{12}\gamma$ en s$^{-1}$.

Pour $\sigma=0.1$ THz, les douze composantes sont reconstruites avec un
écart absolu maximal de $6.94\times10^{-18}$ THz.
Cet écart divisé par la plus grande demi-largeur vaut $2.96\times10^{-16}$ ;
ce n'est pas l'erreur relative maximale branche par branche.
Le programme n'appelle pas le noyau d'accumulation du logiciel.
Il réutilise ses interactions, phonons et constantes : l'indépendance
porte sur la somme, pas sur toute la chaîne microscopique.

\section{Contrôles positifs et échec conservé}
Les demi-largeurs de deux points réciproques concordent à moins de
$2\times10^{-13}$ après normalisation par la composante maximale.
Mais les rapports à la valeur obtenue avec $\sigma=0.1$ THz couvrent
$0.463$ à $1.663$ pour $\sigma=0.05$ THz, et $0.707$ à $2.933$ pour
$\sigma=0.2$ THz. La grille ne fournit donc pas des taux convergés.
Ces variations ne constituent pas un intervalle d'incertitude physique.

Le contrôle des représentants équivalents de frontière échoue :
l'écart atteint $6.8181\times10^{-8}$ THz, pour un seuil
déclaré de $10^{-10}$ THz. Ce seuil reste inchangé.
L'écart apparaît avec la correction polaire dans les moteurs C et
Rust. Sans cette correction, il retombe près de l'arrondi.
Les résidus des paires propres sont de l'ordre de $10^{-15}$ :
une diagonalisation précise ne suffit pas à valider la périodicité.

\section{Expérience discriminante sur la correction polaire}
Pour les représentants $(-1,-1,1)/3$ et $(-1,2,1)/3$, les matrices sont
comparées après la transformation de phase atomique appropriée.
On augmente uniquement la coupure réciproque, en maintenant
fixe le paramètre de séparation $\Lambda=0.2661577394$ :
\begin{center}
\begin{tabular}{ccc}
\toprule
Facteur de coupure & Nombre de vecteurs & Défaut matriciel relatif\\
\midrule
1.00 & 311 & $1.35\times10^{-8}$\\
1.25 & 579 & $1.55\times10^{-12}$\\
1.50 & 1013 & $4.94\times10^{-16}$\\
2.00 & 2381 & $4.94\times10^{-16}$\\
\bottomrule
\end{tabular}
\end{center}
Ce contrôle, vérifié indépendamment sur les matrices sauvegardées,
soutient une explication par troncature pour cette paire.
Il ne corrige pas rétroactivement l'export initial.

\subsection{Continuation : export recalculé}
Les exports complets au point externe sélectionné ont ensuite été
recalculés avec les facteurs de coupure 1, 1.5 et 2, à $\Lambda$ fixé.
Les facteurs 1.5 et 2 passent le seuil initial sur les 18 paires de
représentants dupliqués : l'écart maximal vaut $2.49\times10^{-14}$ THz.
L'inspection initiale et la reconstruction indépendante passent sans
modification de leurs critères.
Les demi-largeurs changent d'au plus $4.63\times10^{-7}$ relativement
au calcul initial, branche par branche. Entre les facteurs 1.5 et 2,
cet écart relatif maximal est $1.39\times10^{-13}$.

Cependant, les éléments individuels de \texttt{pp} diffèrent encore
de $0.110$ après normalisation par le plus grand élément.
La stabilité des demi-largeurs ne certifie donc pas une correspondance
élément par élément des interactions. Le contrôle indépendant montre que les 4296 éléments modifiés au-delà
de $10^{-10}$ du maximum impliquent au moins une branche dégénérée.
Les sommes sur les blocs dégénérés complets concordent à
$6.29\times10^{-15}$ après normalisation globale.
Les projecteurs des sous-espaces restent stables alors que leurs bases
changent : ces observations soutiennent une redistribution entre
vecteurs propres dégénérés. Elles ne reconstruisent pas les phases
complexes des amplitudes et ne prouvent pas l'équivalence des opérateurs.

\section{Ce qui reste à établir}
Quatre branches indépendantes puis quatre revues adversariales séparent
l'accord scalaire du problème d'opérateur. Un facteur amont commun erroné
pourrait préserver cet accord tout en changeant tous les taux absolus.
Des expansions différentes peuvent préserver une diagonale
sans préserver les termes hors diagonale.

Il faut vérifier l'export complet avec correction polaire contrôlée,
la normalisation microscopique indépendante, les permutations de modes,
le comptage des canaux et une intégration conservant les invariants.
Une contribution gaussienne désaccordée n'est pas un événement exactement
résonant de la note 19. La trajectoire en température, les frontières et
la comparaison expérimentale restent ouvertes.

\section{Traçabilité}
Les scripts, versions, empreintes, résultats compacts, contrôles échoués
et rapports de revue sont conservés dans le dossier de campagne.
Les données volumineuses restent sur D. Le programme
\texttt{scripts/run\_aln\_interaction\_pilot.py} reproduit le calcul
dans un répertoire neuf et enregistre explicitement l'échec initial.

La formulation suit Togo, Chaput et Tanaka,
Phys.\ Rev.\ B \textbf{91}, 094306 (2015), ainsi que la documentation
et le code source épinglé de phono3py. Aucune nouveauté théorique,
validation formelle ou classification hydrodynamique complète de l'AlN
n'est revendiquée.
\end{document}
